\section*{الفصل \ref{ch:g10:earth-energy-budget} --- ميزانية الطاقة للأرض}

\begin{solution}{exo:g10:earth-energy-budget:1}
\omterm{def:g10:earth-energy-budget:insolation}{الإشماس} في مكان ما هو قدرة ضوء الشمس التي يتلقاها متر مربع أفقي في أعلى الغلاف
الجوي فوقه (متوسطةً على يوم أو على سنة). \omterm{def:g10:earth-energy-budget:albedo}{وألبيدو} سطح ما هو الكسر الذي يعكسه من
ضوء الشمس الساقط عليه. والثلج الحديث (من 0.8 إلى 0.9) يعكس أكثر بكثير من المحيط
المكشوف (0.06).
\end{solution}

\begin{solution}{exo:g10:earth-energy-budget:2}
المنعكس: $100/340 \approx \qty{29}{\percent}$؛ وما يمتصه الهواء:
$80/340 \approx \qty{24}{\percent}$؛ وما يمتصه السطح: $160/340
\approx \qty{47}{\percent}$.
\end{solution}

\begin{solution}{exo:g10:earth-energy-budget:3}
الإسفلت: $800 \times (1 - 0.04) = \qty{768}{W}$. الثلج: $800 \times (1 -
0.85) = \qty{120}{W}$، أي أقل بأكثر من ست مرات.
\end{solution}

\begin{solution}{exo:g10:earth-energy-budget:4}
التبخر: \qty{82}{W/m^2}، مقابل \qty{21}{W/m^2} لتدفئة الهواء. فالماء الذي
يتبخر يأخذ حرارته الكامنة من البحر؛ ويصعد البخار مع الهواء، وحيث يتكاثف قطيرات
سحابية، غالبًا على ارتفاع كيلومترات، يعيد تلك الطاقة إلى الهواء المحيط به.
\end{solution}

\begin{solution}{exo:g10:earth-energy-budget:5}
في ديسمبر يميل نصف الكرة الشمالي بعيدًا عن الشمس: فتبقى الشمس عند القطب تحت الأفق
طوال اليوم، فلا يصل أي ضوء شمس. وفي الانقلاب الشتوي في ديسمبر يصح ذلك في كل
مكان وراء $90^\circ - 23.44^\circ =
\qty{66.56}{\degree}$ ش.
\end{solution}

\begin{solution}{exo:g10:earth-energy-budget:6}
$1361 \sin 90^\circ = \qty{1361}{W}$؛ $1361 \sin 30^\circ =
\qty{680.5}{W}$؛ $1361 \sin 10^\circ \approx \qty{236}{W}$. ونصف $S_0$
يتطلب $\sin h = 0.5$، أي $h = \qty{30}{\degree}$.
\end{solution}

\begin{solution}{exo:g10:earth-energy-budget:7}
المنعكس 100، وما يمتصه الهواء 80، وما يمتصه السطح 160:
$100 + 80 + 160 = 340$، فكل ضوء الشمس محسوب. وعند السطح،
$185 = 25 + 160$: فما لا ينعكس يُمتص.
\end{solution}

\begin{solution}{exo:g10:earth-energy-budget:8}
القرص: $\pi \times (6.371\times10^6)^2 \approx \qty{1.275e14}{m^2}$؛ والقدرة:
$1361 \times 1.275\times10^{14} \approx \qty{1.74e17}{W}$. والكرة: $4\pi
R^2 \approx \qty{5.10e14}{m^2}$. والمتوسط: $1.74\times10^{17}/5.10\times
10^{14} \approx \qty{340}{W/m^2}$، أي ربع $S_0$.
\end{solution}

\begin{solution}{exo:g10:earth-energy-budget:9}
$0.7 \times 5.10\times10^{14} \times 3.16\times10^7 \approx
\qty{1.1e22}{J}$ في السنة.
\end{solution}

\begin{solution}{exo:g10:earth-energy-budget:10}
تعكس السحب من ضوء الشمس $100 - 53 = \qty{47}{W/m^2}$ إضافية، لا تُمتص إذن.
وتحبس $267 - 239 = \qty{28}{W/m^2}$ من الأشعة تحت الحمراء. وفي المحصلة تفقد
الأرض $47 - 28 = \qty{19}{W/m^2}$: فالسحب، في مجملها، تبردها.
\end{solution}

\begin{solution}{exo:g10:earth-energy-budget:11}
من نحو 10 مايو إلى نهاية يوليو، أي في مايو ويونيو ويوليو. فعلى خط الاستواء تمر
الشمس فوق الرأس عند الظهيرة في الاعتدالين، في مارس وسبتمبر؛ أما في الانقلابين
فتمر على بعد \qty{23.44}{\degree} من سمت الرأس، فيتلقى خط الاستواء أقل قليلًا في
يونيو وفي ديسمبر.
\end{solution}

\begin{solution}{exo:g10:earth-energy-budget:12}
في كل ثانية: $82/(2.47\times10^6) \approx \qty{3.32e-5}{kg}$. وفي كل يوم:
$3.32\times10^{-5} \times 86\,400 \approx \qty{2.9}{kg}$، أي طبقة سماكتها نحو
\qty{2.9}{mm}. وفي كل سنة: $3.32\times10^{-5} \times 3.16\times10^7 \approx
\qty{1050}{kg}$، أي نحو \qty{1.05}{m} من الماء.
\end{solution}

\begin{solution}{exo:g10:earth-energy-budget:13}
$S_0 \sin\delta = S_0/\pi$ يعطي $\sin\delta = 1/\pi \approx 0.318$، إذن
$\delta \approx \qty{18.6}{\degree}$. ثم $\sin\lambda = 0.318/\sin 23.44^\circ
\approx 0.318/0.398 \approx 0.800$: $\lambda \approx \qty{53.15}{\degree}$ أو
$180^\circ - 53.15^\circ = \qty{126.85}{\degree}$. وبين القيمتين تزداد $\lambda$
بمقدار \qty{73.7}{\degree}، أي $73.7/360$ من السنة، نحو
\num{75} يومًا، من منتصف مايو إلى أواخر يوليو.
\end{solution}

\begin{solution}{exo:g10:earth-energy-budget:14}
يمتص الجليد $1 - 0.6 = 0.40$ من ضوء الشمس، والمحيط $1 - 0.06 =
0.94$. القدرة الإضافية: $(0.94 - 0.40) \times 300 \times 10^{12} \approx
\qty{1.6e14}{W}$، أي أكثر من التركيب الضوئي كله على الكوكب. وهذه الطاقة الإضافية
تدفئ الماء، فيذيب مزيدًا من الجليد عند حواف الطوافات، فيتكشف مزيد من الماء
القاتم: فالتغير يغذي نفسه.
\end{solution}

\begin{solution}{exo:g10:earth-energy-budget:15}
في أقرب نقطة: $1361 \times (149.6/147.1)^2 \approx \qty{1408}{W/m^2}$؛ وفي أبعد
نقطة: $1361 \times (149.6/152.1)^2 \approx \qty{1317}{W/m^2}$؛ والفرق بينهما نحو
\qty{7}{\percent}. أما عند \qty{60}{\degree} ش \omterm{def:g10:earth-energy-budget:insolation}{فالإشماس} في يونيو أكبر منه في
ديسمبر بعشرين مرة: ولا يمكن لتغير ببضعة في المئة أن يسببه. ثم إن الأرض تكون أقرب
ما تكون إلى الشمس في أوائل يناير، في أثناء الشتاء الشمالي، ونصفا الكرة، على
المسافة نفسها من الشمس، لهما فصلان متعاكسان: فميل المحور، لا المسافة، هو الذي
يصنع الفصول.
\end{solution}

\begin{solution}{pb:g10:earth-energy-budget:1}
\textbf{1.} $\pi \times (6.371\times10^6)^2 \approx \qty{1.275e14}{m^2}$.

\textbf{2.} $1361 \times 1.275\times10^{14} \approx \qty{1.735e17}{W}$.

\textbf{3.} $4\pi R^2 \approx \qty{5.101e14}{m^2}$.

\textbf{4.} $1.735\times10^{17}/5.101\times10^{14} \approx \qty{340}{W/m^2}$،
أي $S_0/4$.

\textbf{5.} خط الاستواء: \qty{1361}{W}؛ و\qty{60}{\degree} ش: $1361 \sin
30^\circ = \qty{680.5}{W}$.

\textbf{6.} $100/340 \approx 0.294$.

\textbf{7.} $340 - 100 = \qty{240}{W/m^2}$؛ وعلى الأرض كلها، $240 \times
5.101\times10^{14} \approx \qty{1.224e17}{W}$.

\textbf{8.} $53/340 \approx 0.156$. فالسحب ترفع \omterm{def:g10:earth-energy-budget:albedo}{الألبيدو} بعامل
$100/53 \approx 1.9$: إنها تضاعفه تقريبًا.

\textbf{9.} $(240/5.67\times10^{-8})^{1/4} \approx \qty{255}{K}$، أي أدنى بنحو
\qty{33}{K} من درجة حرارة السطح \qty{288}{K}.

\textbf{10.} $160 + 342 = \qty{502}{W/m^2}$.

\textbf{11.} $82 + 21 + 398 = \qty{501}{W/m^2}$. فالسطح يكتسب نحو
\qty{1}{W/m^2}: وهو اختلال التوازن \qty{0.7}{W/m^2} مقرَّبًا.

\textbf{12.} $342/160 \approx 2.1$. فالهواء يرسل إلى السطح على شكل أشعة تحت حمراء
أكثر من ضعف ضوء الشمس الذي يمتصه: إنه أثر الدفيئة.

\textbf{13.} $82 \times 3.156\times10^7/(2.47\times10^6) \approx
\qty{1050}{kg}$، أي طبقة سماكتها نحو \qty{1.05}{m}.

\textbf{14.} $0.7 \times 5.101\times10^{14} \approx \qty{3.6e14}{W}$.

\textbf{15.} $2805/72.07 \approx \qty{38.9}{kJ}$ لكل غرام من الكربون.

\textbf{16.} $104.9\times10^9$ طن تساوي $1.049\times10^{17}$~غ؛
$1.049\times10^{17} \times 38.9\times10^3 \approx \qty{4.08e21}{J}$.

\textbf{17.} $4.08\times10^{21}/3.156\times10^7 \approx \qty{1.29e14}{W}$.

\textbf{18.} $1.29\times10^{14}/5.101\times10^{14} \approx
\qty{0.25}{W/m^2}$.

\textbf{19.} $0.7/0.25 \approx 2.8$: فاختلال توازن الأرض يقارب ثلاثة أضعاف
القدرة التي يخزنها التركيب الضوئي.

\textbf{20.} $1.29\times10^{14}/1.224\times10^{17} \approx 0.0011$:
يخزن التركيب الضوئي \textbf{نحو \qty{0.1}{\percent}} من ضوء الشمس الذي تمتصه
الأرض --- أي جزءًا من ألف.
\end{solution}
